\chapter[Schröder--Bernstein i Cantor]{Comparació de conjunts,\texorpdfstring{\\}{ }Schröder--Bernstein i Cantor}\label{cap:cantor}

\begin{objectius}
\apartat{Prerequisits} Del capítol~\ref{cap:parts-interna}: conjunt de parts, pertinença i inclusió internes, el quantificador universal al llarg d'una projecció, i l'ordre \((P(A),\subseteq)\). Del capítol~\ref{cap:imatges}: imatge directa i la seva adjunció. Del capítol~\ref{cap:valors}: complements i definició per casos. Dels capítols~\ref{cap:parts} i~\ref{cap:valors}: que mono i epi es comproven sobre punts globals.
\apartat{Objectius} En acabar el capítol, el lector hauria de poder:
\begin{itemize}
  \item definir la comparació \(A\preceq B\) i l'equipotència com a relacions del metallenguatge, i demostrar que \(\preceq\) és un preordre invariant;
  \item demostrar Schröder--Bernstein construint el menor punt fix d'un operador creixent sobre \(P(A)\) amb el quantificador universal, i explicar per què no hi intervé cap elecció;
  \item demostrar que cap aplicació \(A\to P(A)\) és exhaustiva (argument diagonal) i deduir-ne el teorema de Cantor i la inexistència d'un conjunt universal;
  \item distingir \enquote{no hi ha exhaustiva} de \enquote{no hi ha isomorfisme}, i situar la paradoxa de Russell com el que el llibre evita en no tenir pertinença global.
\end{itemize}
\end{objectius}

Els deu capítols anteriors han construït tota la infraestructura dels
conjunts ---parts, productes i sumes, quocients, imatges, exponencials i
el conjunt de parts, els nombres naturals i la cardinalitat fins al
numerable--- i han comparat conjunts sobretot en el cas finit i en el
numerable. Ara comença la comparació \emph{general} de mides, i el primer
canvi és de punt de vista. Un cardinal no serà, d'entrada, un conjunt
privilegiat que representa una mida, sinó una manera de \emph{comparar}
conjunts: un conjunt és més petit o igual que un altre quan s'hi injecta, i
tots dos tenen la mateixa mida quan són isomorfs. Amb només aquest
vocabulari ---monomorfismes i isomorfismes, res més que el que ja tenim---
s'obtenen els dos teoremes bàsics de la teoria de Cantor. El primer,
\emph{Schröder--Bernstein}, diu que dues injeccions en sentits oposats
basten per construir un isomorfisme: la comparació de mides és
antisimètrica. El segon, el \emph{teorema de Cantor}, diu que el pas al
conjunt de parts sempre produeix una mida estrictament més gran, de manera
que cap univers de conjunts no queda esgotat per una sola mida. Cap dels
dos no necessita cap axioma nou; en particular, cap dels dos no necessita
una successió prèvia d'alephs ni l'elecció. Aquest és el capítol més curt de
la part, i el que fixa el llenguatge de tots els que segueixen.

\section{La comparació de conjunts}\label{sec:comparacio-conjunts}

Recordem que dos conjunts són \emph{equipotents}, \(A\approx B\), quan hi ha
un isomorfisme \(A\cong B\) (definició~\ref{def:equipotencia}). L'equipotència
és la igualtat de mida. La comparació que falta ---quan una mida és
\emph{menor o igual} que una altra--- s'expressa amb un monomorfisme.

\begin{definicio}[Comparació de conjunts]\index{Comparació de conjunts}\label{def:comparacio}
Siguin \(A\) i \(B\) conjunts.
\begin{itemize}
\item \(A\preceq B\) (\emph{\(A\) es domina per \(B\)}) si existeix un
monomorfisme \(A\hookrightarrow B\).
\item \(A\approx B\) (\emph{\(A\) és equipotent a \(B\)}) si existeix un
isomorfisme \(A\cong B\) (definició~\ref{def:equipotencia}).
\item \(A\prec B\) (\emph{\(A\) és estrictament menor que \(B\)}) si
\(A\preceq B\) però \(A\not\approx B\).
\end{itemize}
\end{definicio}

\begin{observacio}[Una comparació entre conjunts, no una relació interna]\label{obs:comparacio-metateoria}
Com l'equipotència (definició~\ref{def:equipotencia}), les relacions
\(\preceq\), \(\approx\) i \(\prec\) són afirmacions \emph{sobre} conjunts,
formulades al metallenguatge: diuen que existeix una certa aplicació entre
dos conjunts donats. No són parts internes d'un producte ni relacions sobre
un suposat «conjunt de tots els conjunts», que no existeix
(secció~\ref{sec:cantor}). Aquesta distinció, ja feta per a \(\approx\), és
la que permet parlar de mides sense convertir la pertinença global en una
noció primitiva; la mantindrem al llarg de tota la part transfinita.
\end{observacio}

La comparació és un preordre ---reflexiva i transitiva--- i respecta
l'equipotència. Són fets immediats, però convé enunciar-los perquè fixen que
\(\preceq\) mesura una mida i no una còpia concreta.

\begin{proposicio}[La comparació és un preordre invariant]\label{prop:preceq-preordre}
Per a conjunts qualssevol:
\begin{enumerate}
\item \(A\preceq A\) (reflexivitat).
\item Si \(A\preceq B\) i \(B\preceq C\), aleshores \(A\preceq C\)
(transitivitat).
\item Si \(A\approx A'\), \(B\approx B'\) i \(A\preceq B\), aleshores
\(A'\preceq B'\) (invariància per equipotència).
\item Si \(A\approx B\), aleshores \(A\preceq B\) i \(B\preceq A\).
\end{enumerate}
Cap enunciat no usa cap axioma més enllà dels de composició.
\end{proposicio}

\begin{prova}
(1) La identitat \(\id_A\colon A\to A\) és un monomorfisme
(definició~\ref{def:monomorfisme}). (2) La composició de dos monomorfismes és
un monomorfisme: si \(m\colon A\hookrightarrow B\) i \(n\colon
B\hookrightarrow C\) són monomorfismes i \((n\circ m)\circ f=(n\circ m)\circ
g\), aleshores \(n\circ(m\circ f)=n\circ(m\circ g)\) dona \(m\circ f=m\circ
g\) i, altre cop, \(f=g\). (3) Siguin \(u\colon A'\cong A\), \(v\colon B\cong
B'\) isomorfismes i \(m\colon A\hookrightarrow B\) un monomorfisme; aleshores
\(v\circ m\circ u\colon A'\to B'\) és un monomorfisme, perquè els
isomorfismes són monomorfismes (composició amb la inversa) i la classe dels
monomorfismes és tancada per composició. (4) Un isomorfisme és alhora un
monomorfisme i té una inversa, que també ho és; d'aquí \(A\preceq B\) i
\(B\preceq A\).
\end{prova}

L'enunciat que \emph{falta} a la proposició anterior és el recíproc del
punt~(4): si \(A\preceq B\) i \(B\preceq A\), és \(A\approx B\)? És a dir, la
comparació és \emph{antisimètrica} llevat d'equipotència? Que la resposta
sigui afirmativa ---i sense cap axioma d'elecció--- és el teorema de
Schröder--Bernstein, l'objecte de la secció següent.

\section{El teorema de Schröder--Bernstein}\label{sec:schroder-bernstein}

Dues injeccions en sentits oposats, \(f\colon A\hookrightarrow B\) i
\(g\colon B\hookrightarrow A\), no donen directament cap isomorfisme: no hi
ha cap raó perquè \(f\) sigui exhaustiva, ni \(g\). El que sí que es pot fer
és \emph{combinar-les} en una sola bijecció, decidint per a cada element
d'\(A\) si l'enviem per \(f\) o pel tros invertible de \(g\). La decisió es
pren identificant una part distingida d'\(A\) ---els «punts bons»--- sobre la
qual cal usar \(f\), i el seu complement, sobre el qual cal usar
\(g^{-1}\). Tota la construcció viu dins \(P(A)\) i no fa servir cap tria.

\begin{teorema}[Schröder--Bernstein]\index{Schröder--Bernstein, teorema de}\label{teo:schroder-bernstein}
Siguin \(A\) i \(B\) conjunts. Si \(A\preceq B\) i \(B\preceq A\), aleshores
\(A\approx B\). És a dir, de dos monomorfismes \(f\colon A\hookrightarrow B\)
i \(g\colon B\hookrightarrow A\) se'n construeix un isomorfisme
\(A\cong B\).\par
\emph{Axiomes usats:} conjunt de parts, inclusió i pertinença internes i el
quantificador universal \(\forall\) (capítols~\ref{cap:exponencials}
i~\ref{cap:parts-interna}), complement, implicació booleana i definició per
casos (capítol~\ref{cap:valors}), imatge directa i la seva adjunció
(capítol~\ref{cap:imatges}). Cap axioma nou; en particular, \emph{cap
elecció}.
\end{teorema}

\begin{prova}
Escrivim \(g(Y)\hookrightarrow A\) per a la imatge directa
(definició~\ref{def:imatge-directa}) d'una part \(Y\hookrightarrow B\), i
igualment \(f(X)\hookrightarrow B\) per a \(X\hookrightarrow A\). Considerem
l'operació sobre parts d'\(A\)
\[
\Phi\colon P(A)\longrightarrow P(A),
\qquad
\Phi(X)\ :=\ A\setminus g\bigl(B\setminus f(X)\bigr),
\]
composició d'imatges directes i complements
(capítols~\ref{cap:imatges} i~\ref{cap:valors}); és una aplicació interna
\(P(A)\to P(A)\) ben definida. Observem que \(\Phi\) és \emph{creixent} per
la inclusió: si \(X\subseteq X'\), aleshores \(f(X)\subseteq f(X')\), perquè
la imatge directa és monòtona ---és adjunta per l'esquerra de la imatge
inversa (proposició~\ref{prop:adjuncio-imatges})---; d'on \(B\setminus
f(X')\subseteq B\setminus f(X)\), perquè el complement inverteix la inclusió
(capítol~\ref{cap:valors}); d'on \(g(B\setminus f(X'))\subseteq g(B\setminus
f(X))\) altre cop per monotonia de \(g_{*}\); i finalment \(\Phi(X)\subseteq
\Phi(X')\) prenent complements. El nombre parell d'inversions del complement
retorna el sentit de la inclusió.

\emph{Un punt fix de \(\Phi\).} Diem que una part \(X\hookrightarrow A\) és
\emph{\(\Phi\)-tancada} si \(\Phi(X)\subseteq X\). Volem la intersecció de
totes les parts \(\Phi\)-tancades. Convé construir-la amb cura, perquè és un
ínfim \emph{arbitrari} ---indexat per una part de \(P(A)\)--- i el llibre
només té ínfims binaris directament; l'obtenim, però, dins conjunts ja
construïts, mitjançant el quantificador universal. Considerem sobre el
producte \(A\times P(A)\) el predicat intern
\[
\varphi(a,X)\ :\equiv\ \bigl(\Phi(X)\subseteq X\bigr)\Rightarrow
\bigl(a\mathrel{\varepsilon} X\bigr),
\]
on l'antecedent \(\Phi(X)\subseteq X\) és la inclusió interna
(definició~\ref{def:inclusio-interna}), avaluada component \(\Phi\) amb la
projecció, el consequent \(a\mathrel{\varepsilon} X\) és la pertinença
interna (definició~\ref{def:pertinenca-interna}), i la implicació és la
booleana interna del capítol~\ref{cap:valors}. Sigui
\(R_\varphi\hookrightarrow A\times P(A)\) la part que classifica \(\varphi\).
Aplicant a \(R_\varphi\) el quantificador universal \(\forall_{\pi}\) associat
a la projecció \(\pi\colon A\times P(A)\to A\)
(definició~\ref{def:quantificadors}), s'obté una part
\[
D\ :=\ \forall_{\pi}(R_\varphi)\ \hookrightarrow\ A,
\]
els elements \(a\) de la qual són, per la lectura de \(\forall_{\pi}\), els
que compleixen \(\varphi(a,X)\) \emph{per a tota} \(X\); és a dir, els \(a\)
que pertanyen a tota part \(\Phi\)-tancada. Escriurem informalment
\(D=\bigcap\{X\in P(A):\Phi(X)\subseteq X\}\), amb el benentès que aquesta
intersecció arbitrària és, per definició, la part \(\forall_{\pi}(R_\varphi)\)
tot just construïda dins \(A\). La col·lecció de parts \(\Phi\)-tancades és
no buida: la part total \(A\) ho és, perquè \(\Phi(A)\subseteq A\) sempre.
Vegem que \(\Phi(D)=D\).

Una precisió abans de seguir. Per la lectura de \(\forall_\pi\)
(proposició~\ref{prop:quantificadors-lectura}(3)), les \(X\) que intervenen a
la definició de \(D\) no són només parts globals d'\(A\), sinó elements
generalitzats \(X\colon Y\to P(A)\), és a dir, parts \(s_X\hookrightarrow
Y\times A\) a un estadi \(Y\). Els arguments següents valen a qualsevol estadi:
\(\Phi\) és una aplicació interna, composició de \(\exists_f\), \(\neg_B\),
\(\exists_g\) i \(\neg_A\), i sobre les parts associades actua per imatges
directes al llarg de \(\id_Y\times f\) i \(\id_Y\times g\) i per complements
(proposicions~\ref{prop:quantificadors-lectura}(1)
i~\ref{prop:booleanes-puntuals}). Per tant la monotonia que acabem de
provar val igualment a cada estadi, i podem raonar amb \(X\) arbitràries.
Primer, \(\Phi(D)\subseteq D\). Sigui \(X\) qualsevol part
\(\Phi\)-tancada. Com que \(D\subseteq X\) i \(\Phi\) és creixent,
\(\Phi(D)\subseteq\Phi(X)\subseteq X\). Així \(\Phi(D)\subseteq X\) per a
tota \(X\) \(\Phi\)-tancada, d'on \(\Phi(D)\subseteq D\), que és la
intersecció de totes. Per tant \(D\) és ella mateixa \(\Phi\)-tancada.
Segon, \(D\subseteq\Phi(D)\). Com que \(\Phi(D)\subseteq D\) i \(\Phi\) és
creixent, \(\Phi(\Phi(D))\subseteq\Phi(D)\); és a dir, \(\Phi(D)\) és també
\(\Phi\)-tancada. Per la definició de \(D\) com a ínfim, \(D\subseteq
\Phi(D)\). Les dues inclusions donen \(\Phi(D)=D\).

\emph{L'isomorfisme.} Despleguem la igualtat \(D=\Phi(D)=A\setminus
g(B\setminus f(D))\). Prenent complements a \(A\),
\[
A\setminus D\ =\ g\bigl(B\setminus f(D)\bigr).
\]
Aquesta és la clau: sobre \(D\) usarem \(f\), i sobre el complement
\(A\setminus D\) usarem la inversa de \(g\). Definim
\(h\colon A\to B\) per casos (capítol~\ref{cap:valors}), segons la
part \(D\) i el seu complement (teorema~\ref{teo:part-complementari}):
\[
h(a)\ :=\
\begin{cases}
f(a) & \text{si } a\in D,\\[1mm]
g^{-1}(a) & \text{si } a\in A\setminus D.
\end{cases}
\]
La segona clàusula té sentit, i és una aplicació genuïna: la
restricció de \(g\) a \(B\setminus f(D)\) és un monomorfisme, i tot
monomorfisme és un isomorfisme sobre la seva imatge (la seva factorització
imatge té el factor \(e\) isomorfisme, proposició~\ref{prop:mono-nucli-trivial});
la imatge és \(g(B\setminus f(D))=A\setminus D\), de manera que hi ha un
isomorfisme \(A\setminus D\cong B\setminus f(D)\), que és el que escrivim
\(g^{-1}\). Així \(h=[f|_D,\,g^{-1}]\colon D+(A\setminus D)\cong A\to B\) és la
definició per casos sobre la descomposició de \(A\)
(teorema~\ref{teo:part-complementari}).

Comprovem que \(h\) és un isomorfisme veient que és alhora un monomorfisme
i un epimorfisme. Per la separació global, n'hi ha prou amb comprovar-ho
sobre \emph{punts globals}: una aplicació injectiva sobre punts és mono
(observació~\ref{obs:injectiva-mono}) i una d'exhaustiva sobre punts és
epi (observació~\ref{obs:exhaustiva-epi}); i sobre punts, pertànyer a
\(D\) o al seu complement, i a \(f(D)\) o al seu complement, són
alternatives excloents (corol·lari~\ref{cor:elements-de-2}). En el que
segueix, doncs, \(a,a',b\) són punts.

\emph{Injectivitat.} Suposem \(h(a)=h(a')\). Si tots dos són a \(D\),
\(f(a)=f(a')\) dona \(a=a'\) perquè \(f\) és un monomorfisme. Si tots dos són
a \(A\setminus D\), \(g^{-1}(a)=g^{-1}(a')\) dona \(a=g(g^{-1}(a))=
g(g^{-1}(a'))=a'\). Queda descartar el cas mixt: \(a\in D\) i \(a'\in
A\setminus D\). Aleshores \(h(a)=f(a)\in f(D)\), mentre que
\(h(a')=g^{-1}(a')\) amb \(a'=g(b')\) i \(b'\in B\setminus f(D)\); és a dir
\(h(a')=b'\notin f(D)\). Per tant \(h(a)\in f(D)\) i \(h(a')\notin f(D)\), de
manera que \(h(a)\neq h(a')\), contra la hipòtesi. El cas mixt és, doncs,
impossible, i \(h\) és un monomorfisme.

\emph{Exhaustivitat.} Sigui \(b\) un punt de \(B\). Distingim si
\(b\in f(D)\) o no. Si \(b\in f(D)\), aleshores \(b=f(a)\) amb
\(a\in D\), i \(h(a)=f(a)=b\). Si \(b\in B\setminus f(D)\), aleshores
\(g(b)\in g(B\setminus f(D))=A\setminus D\), de manera que
\(h(g(b))=g^{-1}(g(b))=b\). En tots dos casos \(b\) té preimatge per \(h\);
\(h\) és exhaustiva sobre punts i, doncs, un epimorfisme.

Un bimorfisme és un isomorfisme
(proposició~\ref{prop:bimorfisme-isomorfisme}); per tant \(h\colon A\cong
B\) i \(A\approx B\).
\end{prova}

\begin{observacio}[Per què no cal elecció]\label{obs:sb-sense-eleccio}
La construcció de \(D\) no tria res: \(D\) és l'\emph{ínfim} d'una família de
parts ---les \(\Phi\)-tancades--- completament determinada per les dades
\(f\) i \(g\), i l'ínfim, quan existeix, és únic i canònic, igual que ho era
el mínim d'una part de \(\N\) (observació~\ref{obs:minim-vs-eleccio}). Tot
viu dins un conjunt de parts \(P(A)\) ja construït, i la partició d'\(A\) en
\(D\) i el seu complement és una decisió \emph{booleana} interna, no una
successió de tries. És el mateix fenomen del capítol anterior: quan hi ha un
element distingit canònic ---aquí, l'ínfim--- l'elecció és innecessària. La
prova alternativa, que itera \(A_0=A\setminus g(B)\), \(A_{n+1}=g(f(A_n))\) i
en pren la unió, també és correcta, però demana reunir una família indexada
per \(\N\) i comprovar que està uniformement definida; la via de l'ínfim
evita aquesta reunió i és, en aquest llibre, la més neta.
\end{observacio}

\begin{corollari}[Antisimetria de la comparació]\label{cor:preceq-antisimetric}
Per a conjunts qualssevol, \(A\preceq B\) i \(B\preceq A\) impliquen
\(A\approx B\). Per tant la comparació \(\preceq\) és un ordre parcial
\emph{llevat d'equipotència}: reflexiva, transitiva
(proposició~\ref{prop:preceq-preordre}) i antisimètrica mòdul \(\approx\).
\end{corollari}

\begin{prova}
És exactament el teorema~\ref{teo:schroder-bernstein}: de les dues
dominacions se'n treu un isomorfisme. La reflexivitat i la transitivitat són
la proposició~\ref{prop:preceq-preordre}.
\end{prova}

\begin{observacio}[Lectura clàssica]\label{obs:sb-classica}
En termes clàssics, Schröder--Bernstein diu que si hi ha una injecció de
\(A\) en \(B\) i una de \(B\) en \(A\), aleshores hi ha una bijecció entre
\(A\) i \(B\): \(\abs{A}\le\abs{B}\) i \(\abs{B}\le\abs{A}\) donen
\(\abs{A}=\abs{B}\). És el que fa que els cardinals estiguin \emph{totalment
ordenats} un cop se sap, a més, que dos conjunts qualssevol són comparables
---cosa que, aquella sí, requereix l'elecció i que veurem més endavant. Aquí
només hem establert l'antisimetria, que no en necessita cap. La demostració
per punts bons que hem donat és, essencialment, la de Banach; la formulació
amb l'operador creixent \(\Phi\) i el seu menor punt fix és el nucli del
teorema de Knaster--Tarski sobre reticles complets, del qual \(P(A)\) n'és un
exemple (teorema~\ref{teo:inclusio-ordre}).
\end{observacio}

\section{L'argument diagonal i el teorema de Cantor}\label{sec:cantor}

El segon teorema compara un conjunt amb el seu conjunt de parts. La
desigualtat fàcil ---que \(A\) s'injecta en \(P(A)\)--- ja fixa que
\(A\preceq P(A)\). El contingut és que la desigualtat és \emph{estricta}: no
hi pot haver cap isomorfisme \(A\cong P(A)\), ni tan sols cap aplicació
exhaustiva \(A\to P(A)\). La raó és un únic argument, la
\emph{diagonal} de Cantor, que treballa amb la pertinença interna del
capítol~\ref{cap:parts-interna}.

\begin{teorema}[Cap aplicació \(A\to P(A)\) no és exhaustiva]\label{teo:diagonal}
Sigui \(A\) un conjunt i \(f\colon A\to P(A)\) una aplicació qualsevol.
Aleshores \(f\) no és exhaustiva. Concretament, la part
\[
D\ :=\ \{\,a\in A\ \mid\ a\not\mathrel{\varepsilon} f(a)\,\}\ \hookrightarrow\ A
\]
---els elements que \emph{no} pertanyen internament a la part que se'ls
assigna--- no és el valor \(f(a)\) de cap \(a\in A\).\par
\emph{Axiomes usats:} conjunt de parts i pertinença interna
(capítol~\ref{cap:parts-interna}), complement (capítol~\ref{cap:valors}).
Cap axioma nou.
\end{teorema}

\begin{prova}
La part \(D\) està ben definida: la pertinença interna \({\in_A}\) és una
part de \(A\times P(A)\) (definició~\ref{def:pertinenca-interna}), i \(D\) és
la imatge inversa per \(\langle\id_A,f\rangle\colon A\to A\times P(A)\) del
seu complement; és a dir, \(D=\{a : \langle a,f(a)\rangle\notin{\in_A}\}\),
tallada amb el complement (capítol~\ref{cap:valors}) i una imatge inversa
(capítol~\ref{cap:parts}). El seu nom
\(\ulcorner D\urcorner\colon\mathbf1\to P(A)\) és un element de \(P(A)\).

Suposem, cap a una contradicció, que \(D\) fos un valor de \(f\): que
existís un element \(d\colon\mathbf1\to A\) amb \(f(d)=\ulcorner
D\urcorner\), és a dir amb \(f(d)=D\) com a parts d'\(A\). Preguntem-nos si
\(d\in D\). Per la definició de \(D\),
\[
d\in D
\quad\Longleftrightarrow\quad
d\not\mathrel{\varepsilon} f(d).
\]
Però \(f(d)=D\), de manera que \(d\mathrel{\varepsilon} f(d)\) equival a
\(d\in D\) (la pertinença interna a una part constant és la pertinença
generalitzada, proposició~\ref{prop:pertinenca-interna-lectura}). Substituint,
\[
d\in D
\quad\Longleftrightarrow\quad
d\notin D,
\]
que és una contradicció: la distinció dels dos valors de veritat
(teorema~\ref{teo:zero-diferent-u}) prohibeix que un element compleixi una
condició si i només si no la compleix. Per tant no existeix cap \(d\) amb
\(f(d)=D\), i \(D\) no és a la imatge de \(f\). Com que \(D\) és un element
de \(P(A)\) que \(f\) no assoleix, \(f\) no és exhaustiva.
\end{prova}

La desigualtat en l'altre sentit és immediata: cada element determina el seu
singletó, i això injecta \(A\) dins \(P(A)\).

\begin{proposicio}[El monomorfisme dels singletons]\label{prop:singletons}
L'aplicació \(\{-\}\colon A\to P(A)\) que envia cada \(a\in A\) al singletó
\(\{a\}\hookrightarrow A\) és un monomorfisme. Per tant \(A\preceq P(A)\).
\end{proposicio}

\begin{prova}
El singletó \(\{a\}\) és la imatge del punt \(a\colon\mathbf1\to A\); en
termes interns, \(\{-\}\colon A\to P(A)\) és la transposada de la
característica de la diagonal \(\Delta\hookrightarrow A\times A\)
(capítols~\ref{cap:relacions} i~\ref{cap:parts-interna}), i és per tant una
aplicació interna ben definida. És un monomorfisme perquè és injectiu
sobre punts (observació~\ref{obs:injectiva-mono}): si \(\{a\}=\{a'\}\) com
a parts d'\(A\) per a dos punts \(a,a'\), com que sempre
\(a\mathrel{\varepsilon}\{a\}\), resulta \(a\mathrel{\varepsilon}\{a'\}\), és a
dir \(a=a'\). Així \(A\preceq P(A)\).
\end{prova}

\begin{teorema}[Cantor]\index{Cantor, teorema de}\label{teo:cantor}
Per a tot conjunt \(A\),
\[
A\ \prec\ P(A).
\]
És a dir, \(A\preceq P(A)\) però \(A\not\approx P(A)\): el conjunt de parts
és estrictament més gran que el conjunt.\par
\emph{Axiomes usats:} els del teorema~\ref{teo:diagonal} i de la
proposició~\ref{prop:singletons}; conjunt de parts, pertinença interna i
complement. Cap axioma nou; cap elecció.
\end{teorema}

\begin{prova}
La dominació \(A\preceq P(A)\) és la proposició~\ref{prop:singletons}. Falta
veure \(A\not\approx P(A)\). Si hi hagués un isomorfisme \(A\cong P(A)\),
seria en particular una aplicació exhaustiva \(A\to P(A)\), cosa impossible
pel teorema~\ref{teo:diagonal}. Per tant no hi ha cap isomorfisme, \(A\not
\approx P(A)\), i combinant amb la dominació, \(A\prec P(A)\).
\end{prova}

\begin{observacio}[No exhaustiva és més fort que no isomorf]\label{obs:cantor-fort}
El teorema~\ref{teo:diagonal} diu una cosa més forta que el teorema de
Cantor: no és només que no hi hagi cap \emph{isomorfisme} \(A\cong P(A)\),
sinó que no hi ha cap aplicació \emph{exhaustiva} \(A\to P(A)\), sigui o no
injectiva. Aquesta és la formulació que convé recordar, perquè és la que es
transporta a altres contextos ---per exemple, per veure que \(\mathbf2^{\N}\)
no és enumerable: una enumeració seria una exhaustiva
\(\N\twoheadrightarrow\mathbf2^{\N}\approx P(\N)\), que el
teorema~\ref{teo:diagonal} prohibeix. Convé mantenir separades, doncs, les
dues negacions: «no hi ha exhaustiva» i «no hi ha monomorfisme» són
afirmacions distintes, i la diagonal dona la primera.
\end{observacio}

\begin{corollari}[No hi ha conjunt universal ni mida màxima]\label{cor:no-universal}
No existeix cap conjunt \(U\) tal que tot conjunt s'injecti en \(U\). En
particular, no hi ha cap conjunt de mida màxima: per a tot \(A\) existeix un
conjunt estrictament més gran, a saber \(P(A)\).
\end{corollari}

\begin{prova}
Si un tal \(U\) existís, aleshores \(P(U)\preceq U\), perquè \(P(U)\) també
seria un conjunt i s'hauria d'injectar en \(U\). Però la
proposició~\ref{prop:singletons} dona \(U\preceq P(U)\), de manera que
Schröder--Bernstein (teorema~\ref{teo:schroder-bernstein}) forçaria
\(U\approx P(U)\), contra el teorema de Cantor
(teorema~\ref{teo:cantor}). Per tant cap conjunt no domina tots els
conjunts. La segona afirmació és el teorema de Cantor mateix:
\(A\prec P(A)\) exhibeix, per a cada \(A\), un conjunt més gran.
\end{prova}

\begin{observacio}[Lectura clàssica i el lloc de la paradoxa]\label{obs:cantor-classica}
Clàssicament, el teorema de Cantor és \(\abs{A}<2^{\abs{A}}\): el cardinal
del conjunt de parts és estrictament més gran que el del conjunt, perquè
\(P(A)\approx\mathbf2^{A}\) (el conjunt de parts és l'exponencial de base
\(\mathbf2\), capítol~\ref{cap:exponencials}). D'aquí surt la torre infinita
de cardinals \(\abs{A}<2^{\abs{A}}<2^{2^{\abs{A}}}<\cdots\), que no s'atura
mai. El corol·lari~\ref{cor:no-universal} és la forma estructural del fet que
«el conjunt de tots els conjunts» no existeix: si existís, contindria el seu
propi conjunt de parts i violaria Cantor. Aquesta és, de fet, la manera
correcta de dir en aquest llibre el que en la teoria clàssica es formula amb
la paradoxa de Russell; aquí no hi ha cap paradoxa perquè mai no hem tractat
la pertinença global entre conjunts com una relació, sinó sempre la
pertinença \emph{a una part d'un conjunt fix}. La diagonal de Cantor
reapareixerà, en versió ordinal, com el teorema de Burali--Forti, i en versió
de famílies, com el teorema de König.
\end{observacio}

\begin{resumcap}
\apartat{Idees centrals}
\begin{enumerate}
  \item Comparar mides és exhibir aplicacions: \(A\preceq B\) és un monomorfisme, \(A\approx B\) un isomorfisme. Són relacions del metallenguatge entre conjunts, no parts d'un producte.
  \item Schröder--Bernstein: dues dominacions oposades donen un isomorfisme. La prova és un menor punt fix d'un operador creixent \(\Phi\) sobre \(P(A)\), obtingut com a ínfim amb el quantificador universal; cap elecció.
  \item L'isomorfisme \(h=[f|_D,g^{-1}]\) es defineix per casos sobre \(D\) i el seu complement, i es comprova que és bimorfisme sobre punts globals, on les alternatives són excloents.
  \item Argument diagonal: cap \(A\to P(A)\) és exhaustiva, perquè la part \(\{a:a\not\mathrel{\varepsilon}f(a)\}\) no és cap valor. Amb els singletons, \(A\prec P(A)\): Cantor.
  \item No hi ha conjunt universal ni mida màxima; la paradoxa de Russell no apareix perquè el llibre no té pertinença global, només pertinença a parts d'un conjunt fix.
\end{enumerate}
\apartat{Errors típics} Creure que dues injeccions oposades donen directament una bijecció (cal construir-la); intentar l'ínfim de les parts \(\Phi\)-tancades sense el quantificador universal (els axiomes només donen ínfims binaris); fer la distinció de casos de la prova de Schröder--Bernstein sobre elements generalitzats, on \enquote{o} no exclou; confondre \enquote{no isomorf} amb \enquote{no exhaustiu} (la diagonal dona la segona, més forta); i llegir el teorema de Cantor com una afirmació sobre un conjunt de tots els cardinals.
\apartat{Lectura recomanada} \cite[§§8.1--8.3]{lawvere-rosebrugh} per a l'argument diagonal en conjunts abstractes i el teorema del punt fix de Lawvere, del qual Cantor és un cas; \cite[sessió~29]{lawvere-schanuel} per a una presentació elemental de la diagonal; \cite[§§22 i 25]{halmos} per a Schröder--Bernstein i el teorema de Cantor en la teoria clàssica; i \cite[cap.~3]{jech2003} per a la demostració clàssica de Schröder--Bernstein i l'aritmètica cardinal que en resulta.
\end{resumcap}
